\documentclass[10pt,a4paper]{article}

%double si on veut une version prof et une version élève. simple si on veut une seule version.
\newcommand{\buildmode}{double}

\InputIfFileExists{version.tex}{}{}
% Version (définie par le script)
\providecommand{\version}{eleve}

\usepackage{feuilleexo}

%%%à modifier à chaque fois

\renewcommand{\niveau}{1G}
\renewcommand{\chapitre}{Suites, généralités}
\renewcommand{\numchap}{1}
\renewcommand{\theme}{Suites}

\begin{document}




%\legendeicones

\vspace{-0.3cm}

\begin{prerequis}
\begin{itemize}
    \item Calculer l'image d'un nombre par une fonction, notation \(f(x)\).
    \item Résoudre une équation ou une inéquation du premier degré.
    \item Placer des points dans un repère à partir d'un tableau de valeurs.
\end{itemize}

\end{prerequis}

\begin{multicols}{2}

\section{Modes de définitions}

\begin{exo}[Une suite définie explicitement][]

On considère la suite \((u_n)\) définie pour tout entier naturel \(n\) par \(u_n = 3n-2\).

\begin{enumerate}[leftmargin=*]
    \item Calculer \(u_0\), \(u_1\),\(u_2\), \(u_3\) et  \(u_5\).
    \item Calculer \(u_{100}\).
\end{enumerate}
\end{exo}

\begin{exo}[Une suite définie par récurrence][]

On considère la suite \((u_n)\) définie pour tout entier naturel \(n\) par \(\left\{\begin{array}{rcl}
u_0 &= & 5  \\
u_{n+1} & = & 4u_n
\end{array}
\right.\).

\begin{enumerate}[leftmargin=*]
    \item Calculer \(u_0\), \(u_1\),\(u_2\), \(u_3\) et  \(u_5\).
    \item Calculer \(u_{100}\).
\end{enumerate}
\end{exo}

\begin{exo}
    Pour chacune des suites suivantes calculer $u_1$, $u_2$ et $u_4$ et préciser leur mode de définition. Justifier vos réponses par une phrase.

    \begin{enumerate}
        \item $u_n=\frac{n}{2}+\frac{3}{4}$
    
        \item $u_0 = 2$ et $u_{n+1}  =  2u_n-3$

        \item $u_n=\frac{n-2}{n^2+2}$

        \item $u_0 = -3$ et $u_{n+1}  =  -3(u_n)^2+2$

    \end{enumerate}
    
\end{exo}

\begin{exo}[Vrai ou faux][appro]
On considère une suite \((u_n)\) définie pour tout entier naturel \(n\). Dire si chacune des affirmation suivante est vraie ou fausse, en justifiant.

\begin{enumerate}[leftmargin=*]
    \item Si la suite \((u_n)\) est définie explicitement, alors on peut calculer \(u_{50}\) sans connaître \(u_{49}\).
    \item Si la suite \(u_n\) est définie par \(u_n = 7n-3\) alors elle est définie par récurrence.
    \item Si on a définie une relation entre \(u_{n+1}\) et \(u_n\), alors on a défini la suite \(u_n\).
    \item Si une suite est définie par récurrence, alors il est nécessaire de calculer \(u_0, u_1, \ldots, u_9\) pour connaître \(u_{10}\).
    \item Si on connait les deux relations suivantes \(u_1 = 0\) et \(u_{n+1} = u_n +4\), alors on a \(u_4=16\)
\end{enumerate}
\end{exo}

\prof{
\begin{itemize}
    \item 1. Vrai : c'est justement l'intérêt d'une formule explicite \(u_n = f(n)\).
    \item 2. Faux : \(u_n = 7n-3\) donne directement \(u_n\) en fonction de \(n\), c'est une définition explicite.
    \item 3. Faux : sans terme initial, la relation de récurrence ne détermine pas une suite unique.
    \item 4. Vrai (en général, sans information supplémentaire) : la relation de récurrence ne donne \(u_{n+1}\) qu'à partir de \(u_n\), il faut donc dérouler le calcul de proche en proche.
    \item 5. Faux : on a \(v_4 = v_3+4 = 0+4 = 4 \neq 16\).
\end{itemize}
}

\begin{exo}[Retrouver une expression explicite][appro]
On donne ci-dessous les sept premiers termes d'une suite \((u_n)\).

\smallskip

\begin{center}
\begin{tabular}{|c|c|c|c|c|c|c|c|}
    \hline
    \(n\) & 0 & 1 & 2 & 3 & 4 & 5 & 6 \\
    \hline
    \(u_n\) & 0 & 3 & 6 & 9 & 12 & 15 & 18 \\
    \hline
\end{tabular}
\end{center}

\begin{enumerate}[leftmargin=*]
    \item Comment passe-t-on de \(n\) à \(u_n\) ?
    \item Proposer une expression explicite \(u_n = f(n)\) cohérente avec ces valeurs.
    \item Vérifier votre formule en calculant \(u_6\) et en la comparant à la valeur du tableau.
\end{enumerate}
\end{exo}

\begin{exo}[Retrouver une définition par récurrence][appro]
On donne ci-dessous les sept premiers termes d'une suite \((v_n)\).

\smallskip

\begin{center}
\begin{tabular}{|c|c|c|c|c|c|c|c|}
    \hline
    \(n\) & 0 & 1 & 2 & 3 & 4 & 5 & 6 \\
    \hline
    \(v_n\) & 1 & 3 & 9 & 27 & 81 & 243 & 729 \\
    \hline
\end{tabular}
\end{center}

\begin{enumerate}[leftmargin=*]
    \item Comment passe-t-on d'un terme au suivant ?
    \item Proposer une définition par récurrence (terme initial et relation entre \(v_{n+1}\) et \(v_n\)) cohérente avec ces valeurs.
    \item Vérifier votre relation en calculant \(v_6\) à partir de \(v_5\).
\end{enumerate}
\end{exo}

\prof{
\begin{itemize}
    \item Ces deux exercices préparent le suivant : ici chaque tableau \og oriente \fg{} naturellement vers un seul mode de définition (arithmétique pour le premier, géométrique pour le second), contrairement au suivant où les deux modes sont demandés sur les mêmes valeurs.
    \item Utile de faire remarquer que dans le premier cas on ajoute toujours \(3\) (indice de définition par récurrence possible aussi : \(u_{n+1}=u_n+3\), avec \(u_0=0\)), et que dans le second on multiplie toujours par \(3\) (\(v_{n+1}=3v_n\), mais une expression explicite \(v_n=3^n\) existe également).
\end{itemize}
}

\begin{exo}[Trouver une suite à partir de ses premiers termes][avance]
On donne les nombres suivants, qui sont les quatre premiers termes d'une suite \((u_n)\) :
\[
u_0 = 1, \quad u_1 = 3, \quad u_2 = 5, \quad u_3 = 7.
\]

\begin{enumerate}[leftmargin=*]
    \item Proposer une expression explicite \(u_n = f(n)\) cohérente avec ces quatre valeurs.
    \item Proposer une définition par récurrence (terme initial et relation entre \(u_{n+1}\) et \(u_n\)) cohérente avec ces mêmes valeurs.
    \item Vérifier que l'expression \(u_n = 2n+1+n(n-1)(n-2)(n-3)\) est elle aussi cohérente avec les quatre valeurs données.
    \item Cette expression est-elle égale à celle trouvée à la question 1 ? Calculer \(u_4\) avec chacune des deux expressions et comparer.
    \item Que peut-on en conclure sur le fait de \og deviner \fg{} une suite à partir de quelques-uns de ses termes seulement ?
\end{enumerate}
\end{exo}

\prof{
\begin{itemize}
    \item Question 5 : conclusion essentielle : quelques termes ne suffisent jamais à prouver qu'une formule explicite est \og la bonne \fg{}. Une formule doit être donnée dans l'énoncé (ou démontrée), pas seulement devinée à partir d'exemples. Cela prépare l'idée de démonstration par récurrence, vue plus tard dans l'année.
     \item Question Bonus : Trouver une suite qui coincide avec \(2n+1\) sur les 8 premiers termes mais pas le 9ème.
     \item 
\end{itemize}
}

\begin{exo}[][]
On considère la suite \((v_n)\) définie pour tout entier naturel \(n\) par \(v_n = n^2-1\).

\begin{enumerate}[leftmargin=*]
    \item Calculer \(v_0\), \(v_1\), \(v_2\) et \(v_3\).
    \item Exprimer, en fonction de \(n\), les termes \(v_{n+1}\) et \(v_{n-1}\).
    \item Le nombre \(48\) est-il un terme de la suite \((v_n)\) ?
\end{enumerate}
\end{exo}

\begin{exo}[][appro]
On considère la suite \((w_n)\) définie pour tout entier naturel \(n\) par \(w_n = \dfrac{2n+1}{n+2}\).

\begin{enumerate}[leftmargin=*]
    \item Justifier que \(w_n\) est bien défini pour tout entier naturel \(n\).
    \item Calculer \(w_0\), \(w_1\) et \(w_2\).
    \item Exprimer \(w_{n+1}\) en fonction de \(n\).
\end{enumerate}
\end{exo}

\prof{
    Objectif : montrer que l'expression de \(u_{n+1}\) en fonction de \(n\) s'obtient en substituant \(n+1\) à \(n\) partout, y compris dans une expression plus complexe (fraction). Anticipe les confusions du type \(w_{n+1} = w_n + 1\).
}

\section{Sens de variation}

\begin{exo}[][]
On considère la suite \((u_n)\) définie pour tout entier naturel \(n\) par \(u_n = 4n-5\).

\begin{enumerate}[leftmargin=*]
    \item Calculer \(u_0\), \(u_1\), \(u_2\) et \(u_3\). Conjecturer le sens de variation de la suite.
    \item Exprimer \(u_{n+1}\)
    \item Calculer \(u_{n+1}-u_n\) en fonction de \(n\).
    \item En déduire le sens de variation de la suite \((u_n)\).
\end{enumerate}
\end{exo}

\prof{
\begin{itemize}
    \item Introduire la méthode générale : étudier le signe de \(u_{n+1}-u_n\).
    \item Faire le lien avec le sens de variation observé numériquement à la question 1, puis justifié algébriquement à la question 2.
    \item Vocabulaire à fixer : croissante, décroissante, constante, et le cas \og non monotone \fg{}.
\end{itemize}
}

\begin{exo}[][]
Pour chacune des suites suivantes, calculer \(u_{n+1}-u_n\) en fonction de \(n\), puis en déduire son sens de variation.

\begin{tasks}(3)
    \task \(u_n = 4n+1\)
    \task \(u_n = -2n+5\)
    \task \(u_n = n^2\)
    \task \(u_n = \dfrac{1}{n+1}\)
    \task \(u_n = 5\)
    \task \(u_n = -n^4\)
\end{tasks}
\end{exo}

\begin{exo}[][appro]
On considère la suite \((u_n)\) définie par récurrence par \(u_0=1\) et, pour tout entier naturel \(n\), \(u_{n+1}=u_n+2n+1\).

\begin{enumerate}[leftmargin=*]
    \item Calculer \(u_1\), \(u_2\) et \(u_3\).
    \item Sans calculer \(u_{n+1}-u_n\) directement, justifier à l'aide de la relation de récurrence que la suite \((u_n)\) est croissante.
\end{enumerate}
\end{exo}

\prof{
    Idée clé : pour une suite définie par récurrence sous la forme \(u_{n+1}=u_n+f(n)\), le signe de \(f(n)\) donne directement le sens de variation, sans calculer explicitement \(u_{n+1}-u_n\). Ici \(2n+1 > 0\) pour tout \(n \geq 0\).
}

\begin{exo}[Vrai ou faux][appro]
Soit \((u_n)\) une suite définie pour tout entier naturel \(n\). Les affirmations suivantes sont-elles vraies ou fausses ? Justifier chaque réponse, à l'aide d'un contre-exemple si besoin.

\begin{enumerate}[leftmargin=*]
    \item Si \(u_0\leq u_1 \leq u_2 \leq u_3\) alors \((u_n)\) est croissante.
    \item Si pour tout entier naturel \(n\), \(u_{n+1}-u_n\) est constant, alors la suite \((u_n)\) est croissante.
    \item Si pour tout entier naturel \(n\), \(u_{n+1}-u_n = 0\) alors la suite \((u_n)\) est constante.
\end{enumerate}
\end{exo}

\prof{

\begin{enumerate}
    \item Nous avons fait attention a bien écrire les phrases sous la forme d'une implication et plus précisément sous la forme "si ... alors".
    \item On a introduit avec les suites le quantificateur universel "pour tout". 
\end{enumerate}
    

}

\section{Représentation graphique : nuage de points}

\begin{exo}[][]
On considère la suite \((u_n)\) définie pour tout entier naturel \(n\) par \(u_n = \frac{1}{2}n^2-2n+3\).

\begin{enumerate}[leftmargin=*]
    \item Calculer \(u_0\), \(u_1\), \(u_2\), \(u_3\) et \(u_4\).
    \item Recopier et compléter le tableau de valeurs suivant.

    \smallskip

    \begin{center}
    \begin{tabular}{|c|c|c|c|c|c|}
        \hline
        \(n\) & 0 & 1 & 2 & 3 & 4 \\
        \hline
        \(u_n\) & & & & & \\
        \hline
    \end{tabular}
    \end{center}

    \item Représenter le nuage de points de coordonnées \((n\,;u_n)\) pour \(n\) allant de \(0\) à \(4\) dans un repère.
    \item À l'aide du nuage de points, conjecturer le sens de variation de la suite \((u_n)\).
\end{enumerate}
\end{exo}

\prof{
\begin{itemize}
    \item Rappeler qu'un nuage de points d'une suite est constitué de points \textbf{isolés}, non reliés entre eux (contrairement à la courbe d'une fonction).
    \item La conjecture graphique de la question 4 sera à confirmer algébriquement (voir section précédente) : bon moment pour articuler les deux approches.
\end{itemize}
}

\begin{exo}[][appro]
On donne ci-dessous le nuage de points représentant les premiers termes d'une suite \((u_n)\).

\begin{center}
\begin{tikzpicture}
\begin{axis}[
    axis lines=middle,
    xmin=-0.5, xmax=6.5,
    ymin=-0.5, ymax=7,
    xtick={0,1,2,3,4,5,6},
    ytick={0,1,2,3,4,5,6},
    xlabel={\(n\)},
    ylabel={\(u_n\)},
    width=8cm,
    height=7cm
]
\addplot[only marks, mark=*, mark size=2pt, blue] coordinates {
    (0,6) (1,4.5) (2,3.5) (3,3) (4,2.8) (5,2.7) (6,2.65)
};
\end{axis}
\end{tikzpicture}
\end{center}

\begin{enumerate}[leftmargin=*]
    \item Donner une valeur approchée de \(u_0\), \(u_2\) et \(u_5\).
    \item Conjecturer le sens de variation de la suite \((u_n)\).
    \item Un élève affirme : \og Comme les points semblent se rapprocher d'une même hauteur, la suite finira par être constante. \fg{} Que penser de cette affirmation ?
\end{enumerate}
\end{exo}

\prof{
    Question 3 : ouverture vers la notion de limite (hors programme à ce stade), à ne pas développer mais à ne pas fermer non plus. L'objectif est de faire percevoir la différence entre \og se rapprocher de \fg{} et \og atteindre \fg{}.
}

\begin{exo}[][appro]
On considère la suite \((u_n)\) définie par récurrence par \(u_0=1\) et, pour tout entier naturel \(n\), \(u_{n+1}=-u_n+4\).

\begin{enumerate}[leftmargin=*]
    \item Calculer \(u_0\), \(u_1\), \(u_2\), \(u_3\) et \(u_4\).
    \item Représenter le nuage de points correspondant.
    \item La suite \((u_n)\) est-elle croissante ? Décroissante ? Que remarque-t-on ?
\end{enumerate}
\end{exo}

\prof{
    Suite oscillante (alternance \(1\); \(3\); \(1\); \(3\); \ldots) : ni croissante, ni décroissante, ni constante. Utile pour montrer que toutes les suites ne sont pas monotones, et que le nuage de points permet de le voir immédiatement.
}

\end{multicols}

\end{document}
}
